% GENERATED FILE. Edit canonical lessons, then run npm run generate:books.
% canonical-source-hash: fnv1a64:181ef0375392c6c3
% canonical-lessons: TA-C430-utayam, TA-C430-veliccam, TA-C430-ativanam, TA-C430-minnu, TA-C430-oli

\chapter{Light in the Sky}
\label{ch:ta-light-in-the-sky}
\hlchaptermodality{430}

\begin{chapteropening}
\textbf{By the end of this chapter:} \emph{I can talk about sunrise, brightness and the horizon, and things glittering and ringing out.}

Complete the last lesson of chapter 430: I can talk about sunrise, brightness and the horizon, and things glittering and ringing out.
\end{chapteropening}

\section[\texorpdfstring{\ta{உதயம்} (utayam)}{உதயம் (utayam)}]{\ta{உதயம்} (utayam) — sunrise; a rising}

\label{lesson:TA-C430-utayam}



\begin{quote}
Before the new one: say the Tamil for a downpour, then the Tamil for cloud cover.
\end{quote}

\subsection*{You'll want to know: \ta{உதயம்}}
\textbf{\ta{உதயம்}} — \emph{utayam} — \textquotedblleft{}sunrise; a rising\textquotedblright{}.

\begin{quote}
\ta{மலை} \ta{உச்சியில்} \ta{சூரிய} \ta{உதயத்தைப்} \ta{பார்த்தோம்}.
\end{quote}

\begin{quote}
\emph{malai ucciyil cūriya utayattaip pārttōm.}
\end{quote}

\begin{quote}
We watched the sunrise from the hilltop.
\end{quote}

\ta{சூரிய} \ta{உதயம்} is sunrise.

\begin{cousinweb}
\textbf{\ta{உதயம்}} (\emph{utayam}) is borrowed from Sanskrit \emph{udaya}, \textquotedblleft{}rise, the rising of the sun\textquotedblright{} (CDIAL 1931).
\end{cousinweb}

\begin{grammarlens}[title={What you've built}]
One more word for getting around.
\end{grammarlens}

\subsection*{Your turn}
\begin{itemize}
\raggedright
  \item \emph{Say it:} \emph{utayam}
  \item \emph{Say it:} \emph{utayam}, once more
  \item \emph{Say it:} \emph{malai ucciyil cūriya utayattaip pārttōm.} — We watched the sunrise from the hilltop.
  \item \emph{Recall:} say the Tamil for a downpour, then the Tamil for cloud cover, then say \emph{utayam} again
\end{itemize}

\begin{tcolorbox}[breakable,colback=teal!4,colframe=teal!35!black,title={Before you move on}]
What does \ta{உதயம்} mean? (\textquotedblleft{}Sunrise; a rising\textquotedblright{}.) Say it once more.
\end{tcolorbox}
\section[\texorpdfstring{\ta{வெளிச்சம்} (veḷiccam)}{வெளிச்சம் (veḷiccam)}]{\ta{வெளிச்சம்} (veḷiccam) — brightness; light}

\label{lesson:TA-C430-veliccam}



\begin{quote}
Before the new one: say the Tamil for cloud cover, then the Tamil for sunrise.
\end{quote}

\subsection*{You'll want to know: \ta{வெளிச்சம்}}
\textbf{\ta{வெளிச்சம்}} — \emph{veḷiccam} — \textquotedblleft{}brightness; light\textquotedblright{}.

\begin{quote}
\ta{சன்னல்} \ta{வழியாக} \ta{வெளிச்சம்} \ta{உள்ளே} \ta{வந்தது}.
\end{quote}

\begin{quote}
\emph{caṉṉal vaḻiyāka veḷiccam uḷḷē vantatu.}
\end{quote}

\begin{quote}
Light came in through the window.
\end{quote}

The opposite is \ta{இருட்டு}, darkness.

\begin{cousinweb}
\textbf{\ta{வெளிச்சம்}} (\emph{veḷiccam}) is native Dravidian (DEDR 5496), \textquotedblleft{}light, clearness\textquotedblright{}, from the root of \ta{வெள்}, \textquotedblleft{}white, bright\textquotedblright{}.
\end{cousinweb}

\begin{grammarlens}[title={What you've built}]
One more word for getting around.
\end{grammarlens}

\subsection*{Your turn}
\begin{itemize}
\raggedright
  \item \emph{Say it:} \emph{veḷiccam}
  \item \emph{Say it:} \emph{veḷiccam}, once more
  \item \emph{Say it:} \emph{caṉṉal vaḻiyāka veḷiccam uḷḷē vantatu.} — Light came in through the window.
  \item \emph{Recall:} say the Tamil for cloud cover, then the Tamil for sunrise, then say \emph{veḷiccam} again
\end{itemize}

\begin{tcolorbox}[breakable,colback=teal!4,colframe=teal!35!black,title={Before you move on}]
What does \ta{வெளிச்சம்} mean? (\textquotedblleft{}Brightness; light\textquotedblright{}.) Say it once more.
\end{tcolorbox}
\section[\texorpdfstring{\ta{அடிவானம்} (aṭivāṉam)}{அடிவானம் (aṭivāṉam)}]{\ta{அடிவானம்} (aṭivāṉam) — the horizon}

\label{lesson:TA-C430-ativanam}



\begin{quote}
Before the new one: say the Tamil for sunrise, then the Tamil for brightness.
\end{quote}

\subsection*{You'll want to know: \ta{அடிவானம்}}
\textbf{\ta{அடிவானம்}} — \emph{aṭivāṉam} — \textquotedblleft{}the horizon\textquotedblright{}.

\begin{quote}
\ta{அடிவானத்தில்} \ta{சூரியன்} \ta{மறைந்தது}.
\end{quote}

\begin{quote}
\emph{aṭivāṉattil cūriyaṉ maṟaintatu.}
\end{quote}

\begin{quote}
The sun disappeared below the horizon.
\end{quote}

\ta{அடி}, the bottom, + \ta{வானம்}, the sky.

\begin{cousinweb}
\textbf{\ta{அடிவானம்}} (\emph{aṭivāṉam}) is built from \ta{அடி}, \textquotedblleft{}the foot, the bottom\textquotedblright{} (DEDR 72), and \ta{வானம்}, \textquotedblleft{}the sky\textquotedblright{} (DEDR 5381).
\end{cousinweb}

\begin{grammarlens}[title={What you've built}]
One more word for getting around.
\end{grammarlens}

\subsection*{Your turn}
\begin{itemize}
\raggedright
  \item \emph{Say it:} \emph{aṭivāṉam}
  \item \emph{Say it:} \emph{aṭivāṉam}, once more
  \item \emph{Say it:} \emph{aṭivāṉattil cūriyaṉ maṟaintatu.} — The sun disappeared below the horizon.
  \item \emph{Recall:} say the Tamil for sunrise, then the Tamil for brightness, then say \emph{aṭivāṉam} again
\end{itemize}

\begin{tcolorbox}[breakable,colback=teal!4,colframe=teal!35!black,title={Before you move on}]
What does \ta{அடிவானம்} mean? (\textquotedblleft{}The horizon\textquotedblright{}.) Say it once more.
\end{tcolorbox}
\section[\texorpdfstring{\ta{மின்னு} (miṉṉu)}{மின்னு (miṉṉu)}]{\ta{மின்னு} (miṉṉu) — to glitter; to flash}

\label{lesson:TA-C430-minnu}



\begin{quote}
Before the new one: say the Tamil for brightness, then the Tamil for the horizon.
\end{quote}

\subsection*{You'll want to know: \ta{மின்னு}}
\textbf{\ta{மின்னு}} — \emph{miṉṉu} — \textquotedblleft{}to glitter; to flash\textquotedblright{}.

\begin{quote}
\ta{இரவில்} \ta{நட்சத்திரங்கள்} \ta{மின்னின}.
\end{quote}

\begin{quote}
\emph{iravil naṭcattiraṅkaḷ miṉṉiṉa.}
\end{quote}

\begin{quote}
At night the stars glittered.
\end{quote}

Past (it …): \ta{மின்னியது}, \emph{miṉṉiyatu}.

\ta{மின்னல்}, lightning, is built on it.

\begin{cousinweb}
\textbf{\ta{மின்னு}} (\emph{miṉṉu}) is native Dravidian (DEDR 4876), \textquotedblleft{}to emit lightning, to shine, to glitter\textquotedblright{}.
\end{cousinweb}

\begin{grammarlens}[title={What you've built}]
One more word for getting around.
\end{grammarlens}

\subsection*{Your turn}
\begin{itemize}
\raggedright
  \item \emph{Say it:} \emph{miṉṉu}
  \item \emph{Say it:} \emph{miṉṉu}, once more
  \item \emph{Say it:} \emph{iravil naṭcattiraṅkaḷ miṉṉiṉa.} — At night the stars glittered.
  \item \emph{Recall:} say the Tamil for brightness, then the Tamil for the horizon, then say \emph{miṉṉu} again
\end{itemize}

\begin{tcolorbox}[breakable,colback=teal!4,colframe=teal!35!black,title={Before you move on}]
What does \ta{மின்னு} mean? (\textquotedblleft{}To glitter; to flash\textquotedblright{}.) Say it once more.
\end{tcolorbox}
\section[\texorpdfstring{\ta{ஒலி} (oli)}{ஒலி (oli)}]{\ta{ஒலி} (oli) — to sound; to ring out}

\label{lesson:TA-C430-oli}



\begin{quote}
Before the new one: say the Tamil for the horizon, then the Tamil for to glitter.
\end{quote}

\subsection*{You'll want to know: \ta{ஒலி}}
\textbf{\ta{ஒலி}} — \emph{oli} — \textquotedblleft{}to sound; to ring out\textquotedblright{}.

\begin{quote}
\ta{ரயில்} \ta{புறப்படும்போது} \ta{மணி} \ta{ஒலித்தது}.
\end{quote}

\begin{quote}
\emph{rayil puṟappaṭumpōtu maṇi olittatu.}
\end{quote}

\begin{quote}
As the train was leaving, a bell rang out.
\end{quote}

Past (it …): \ta{ஒலித்தது}, \emph{olittatu}.

\ta{ஒலி} is also the noun, a sound.

\begin{cousinweb}
\textbf{\ta{ஒலி}} (\emph{oli}) is native Dravidian (DEDR 996), \textquotedblleft{}to sound; a sound, a noise\textquotedblright{}.
\end{cousinweb}

\begin{grammarlens}[title={What you've built}]
One more word for getting around.
\end{grammarlens}

\subsection*{Your turn}
\begin{itemize}
\raggedright
  \item \emph{Say it:} \emph{oli}
  \item \emph{Say it:} \emph{oli}, once more
  \item \emph{Say it:} \emph{rayil puṟappaṭumpōtu maṇi olittatu.} — As the train was leaving, a bell rang out.
  \item \emph{Recall:} say the Tamil for the horizon, then the Tamil for to glitter, then say \emph{oli} again
\end{itemize}

\begin{tcolorbox}[breakable,colback=teal!4,colframe=teal!35!black,title={Before you move on}]
What does \ta{ஒலி} mean? (\textquotedblleft{}To sound; to ring out\textquotedblright{}.) Say it once more.
\end{tcolorbox}
